% !PRINTED 257
\weberpage{257}
%<<CONT
ableiten, wie es durch die Striche angedeutet ist, d.\,h. so, dass
man jede vorhergehende aus der nachfolgenden erhält, indem
man die letzte Zeile und Colonne weglässt; es ist also
\[
R_0 = 1, \quad R_1 = a_{1,1}, \quad R_2 = a_{1,1}a_{2,2} - a_{1,2}^2,
\]
\[
R_3 = \begin{vmatrix} a_{1,1}, & a_{1,2}, & a_{1,3} \\ a_{2,1}, & a_{2,2}, & a_{2,3} \\ a_{3,1}, & a_{3,2}, & a_{3,3} \end{vmatrix}, \ldots R_m = R.
\]
Ein solches System soll, wenn es in absteigender Reihe
\[
R_m,\ R_{m-1} \ldots R_1,\ R_0
\]
geordnet ist, eine Kette von Haupt-Unterdeterminanten heissen.
Solcher Ketten lassen sich $\Pi(m)$ verschiedene bilden, in denen
allen das erste Glied $R$, das letzte Glied 1 ist.

\begin{center}
\S.\,81.

\gesperrt{Quadratische Formen mit verschwindender Determinante.}
\end{center}

Wenn die in (3), \S.\,80 definirte Determinante $R$ der quadra\-tischen
Form $\varphi(x_1, x_2 \ldots x_m)$ verschwindet, so lässt sich, wie wir
schon im \S.\,57 gesehen haben, $\varphi$ durch weniger als $m$ von ein\-ander
unabhängige lineare Functionen von $x$ ausdrücken. Wenn
wir mit $m - \varrho$ die kleinste Zahl von Variablen bezeichnen, durch
die sich $\varphi$ ausdrücken lässt, so hat $\varrho$ dieselbe Bedeutung, wie
im vorigen Paragraphen.

Wenn sich nun unter den Haupt-Unterdeterminanten von $R$
eine $k$-reihige findet, die von Null verschieden ist, etwa
\[
R_k = \Sigma \pm a_{1,1}a_{2,2} \ldots a_{k,k},
\]
so ist, wenn wir
\[
x_{k+1} = 0,\quad x_{k+2} = 0, \ldots x_m = 0
\]
setzen,
\[
\varphi(x_1, x_2 \ldots x_k, 0, 0, 0)
\]
eine quadratische Form von $k$ Variablen $x_1, x_2 \ldots x_k$, deren
Determinante $R_k$ von Null verschieden ist, die sich sonach nicht
durch weniger als $k$ Variable ausdrücken lässt. Um so weniger
kann also bei unbeschränkt veränderlichen $x_1, x_2 \ldots x_m$ die
Function $\varphi$ von weniger als $k$ Variablen abhängen und es folgt
% The source prints the rank/nullity bound with a less-than-or-equal sign.
\weq{1}{\varrho \leqq m - k.}
